\chapter{Results and Evaluation}

\section{Experimental Protocol}
Experiments were conducted on queue-scene videos with different traffic dynamics.
The monitored indicators include counting accuracy, estimation stability, and alert responsiveness.

\section{Observed Metrics}
\begin{itemize}
    \item number of people inside the zone,
    \item arrival rate $\lambda$,
    \item service rate $\mu$,
    \item estimated waiting time $W$,
    \item credible intervals and uncertainty level.
\end{itemize}

\section{Qualitative Analysis}
The system correctly tracks queue-load evolution and flags critical phases.
The addition of credible intervals improves the operational interpretation of results.

\section{Example Results Table}
\begin{table}[H]
    \centering
    \caption{Example experimental summary (to be completed)}
    \begin{tabularx}{\textwidth}{lcccX}
        \toprule
        Scenario & $\lambda$ (pers/s) & $\mu$ (pers/s) & $W$ (s) & Uncertainty Level \\
        \midrule
        Low load & 0.04 & 0.08 & 1.00 & Low \\
        Medium load & 0.08 & 0.11 & 2.67 & Medium \\
        High load & 0.10 & 0.11 & 10.00 & High \\
        \bottomrule
    \end{tabularx}
\end{table}

\section{Evaluation Limitations}
Comprehensive multi-scenario evaluation still needs to be extended (camera-angle variations, strong occlusions, lighting changes, multi-zone setups).
